\chapter{1981:Kenichi Fukui and Roald Hoffmann}

\label{chap:chemistry1981}

The Nobel Prize in Chemistry for the year 1981 was awarded jointly to Kenichi Fukui and Roald Hoffmann.

\textbf{Motivation:}

for their theories, developed independently, concerning the course of chemical reactions

\section{Kenichi Fukui}

\subsection*{Early Life and Education}

Kenichi Fukui was born on 1918-10-04 in Nara, Japan.

\subsection*{Career and Major Research}

Fukui studied at the Kyoto Imperial University, where he received his doctorate and later became professor of chemistry. He developed the frontier orbital theory, which describes how the highest occupied and lowest unoccupied molecular orbitals control the reactivity of molecules.

\subsection*{Nobel Prize-winning Work}

The work for which Kenichi Fukui was awarded the Nobel Prize in Chemistry in 1981 was described by the Nobel Committee as: "for their theories, developed independently, concerning the course of chemical reactions".

Fukui showed that chemical reactions are governed by the interaction of the frontier orbitals, explaining the regio- and stereoselectivity of many reactions.

\subsection*{Significance and Impact}

Frontier orbital theory provided a simple and powerful way to predict the course of chemical reactions and became a standard tool for chemists.

\subsection*{Later Career and Honors}

Fukui remained at Kyoto University until his retirement. He was the first Japanese Nobel laureate in chemistry. He died on 1998-01-09 in Kyoto, Japan.

\subsection*{Legacy}

Fukui's frontier orbital concept underlies modern computational chemistry and the understanding of reactivity.

\subsection*{Selected Publications}

"A Molecular Orbital Theory of Reactivity in Aromatic Hydrocarbons" (Journal of Chemical Physics, 1952)

\clearpage

\section{Roald Hoffmann}

\subsection*{Early Life and Education}

Roald Hoffmann was born on 1937-07-18 in Zloczów, Poland.

\subsection*{Career and Major Research}

Hoffmann survived the Holocaust and emigrated to the United States, studying at Columbia University and Harvard University, where he received his doctorate. He worked at Harvard with Robert B. Woodward and then became professor at Cornell University. He developed extended Huckel calculations and, with Woodward, formulated the Woodward-Hoffmann rules for the conservation of orbital symmetry.

\subsection*{Nobel Prize-winning Work}

The work for which Roald Hoffmann was awarded the Nobel Prize in Chemistry in 1981 was described by the Nobel Committee as: "for their theories, developed independently, concerning the course of chemical reactions".

The Woodward-Hoffmann rules predict the course of pericyclic reactions on the basis of the symmetry of molecular orbitals.

\subsection*{Significance and Impact}

The rules explained how stereochemistry is controlled in pericyclic reactions and helped unify organic reaction theory.

\subsection*{Later Career and Honors}

Hoffmann remained at Cornell University and is also a published writer of poetry and essays. He is still alive.

\subsection*{Legacy}

The Woodward-Hoffmann rules remain a fundamental principle of physical organic chemistry.

\subsection*{Selected Publications}

"The Conservation of Orbital Symmetry" (Angewandte Chemie International Edition, 1969)

\clearpage

\section*{Chapter Summary}

The 1981 prize recognized Fukui and Hoffmann for their independently developed theories of the course of chemical reactions. Their orbital-based theories explained and predicted reactivity in organic chemistry.

